Write \(s_N=s(M_N)\), \(\widehat s_N=\operatorname{vech}(\widehat{\boldsymbol\Sigma})\), \(m=p(p+1)/2\), \(\lambda_1=\lambda\), and \(\lambda_2=1-\lambda\).  All probability statements below are under the triangular sample law generated by \(M_N\).  The local-sequence condition gives \(s_N\to s_0\); since \(s_0\in\Omega_p\), both limiting covariance matrices \(\Sigma_{0,e}\) are positive definite.

For environment \(e\), let \(X_{ei,N}=\operatorname{vech}(Z_{ei}Z_{ei}^{\mathsf T}-\Sigma_{M_N,e})\).  Gaussian fourth moments and Isserlis' identity give
\[
 \operatorname{Cov}(X_{ei,N})_{AB,CD}
 =\Sigma_{M_N,e,AC}\Sigma_{M_N,e,BD}
  +\Sigma_{M_N,e,AD}\Sigma_{M_N,e,BC}
 =V(\Sigma_{M_N,e})_{AB,CD}.
 \tag{1}
\]
Because \(\Sigma_{M_N,e}\to\Sigma_{0,e}\), Gaussian moments of every fixed order are uniformly bounded.  In particular, for every \(\epsilon>0\), Markov's inequality gives
\[
 \mathbb E\!\left[
 \lVert X_{ei,N}\rVert^2
 {\bf1}\{\lVert X_{ei,N}\rVert>\epsilon\sqrt{n_e}\}
 \right]
 \le \frac{\mathbb E\lVert X_{ei,N}\rVert^4}{\epsilon^2n_e}
 \longrightarrow0.
 \tag{2}
\]
Thus the Lindeberg condition in \(\text{lem:multivariate-lindeberg-feller}\) holds.  Moreover, \(\overline Z_e\overline Z_e^{\mathsf T}=O_{\mathbb P_N}(n_e^{-1})\), and replacing \(n_e^{-1}\) by \(\nu_e^{-1}\) also changes the covariance half-vector by \(O_{\mathbb P_N}(n_e^{-1})\).  These are \(o_{\mathbb P_N}(n_e^{-1/2})\) because \(n_e\to\infty\).  Independence of the two samples from \(\text{ass:independent-gaussian-sampling}\), (1)--(2), \(\text{lem:multivariate-lindeberg-feller}\), and \(n_e/N\to\lambda_e\in(0,1)\) from \(\text{ass:asymptotic-sample-balance}\) therefore imply
\[
 \sqrt N(\widehat s_N-s_N)\Rightarrow\mathbb Z_0,
 \qquad
 Y_N:=\sqrt N(\widehat s_N-s_0)
 \Rightarrow y:=h+\mathbb Z_0.
 \tag{3}
\]
The covariance in (3) is \(V_0\) by definition.  It is positive definite: for a nonzero symmetric matrix \(A\), positivity of \(\Sigma_{0,e}\) gives
\[
 \operatorname{Var}(Z_e^{\mathsf T}AZ_e)
 =2\operatorname{tr}\!\left[
   \{\Sigma_{0,e}^{1/2}A\Sigma_{0,e}^{1/2}\}^2
 \right]>0,
 \tag{4}
\]
and both \(\lambda_e\) are positive.  Hence \(y\) has a Lebesgue density on \(\mathbb R^d\).

We next derive the local form of the exact Wishart bodies.  For deterministic \(z=(z_1,z_2)\), let \(\mathcal A_N(z)\) be the set of \(u=(u_1,u_2)\) for which
\[
 S_e(u_e):=\operatorname{vech}^{-1}(s_{0,e}+u_e/\sqrt N)\succ0,
 \qquad
 T_{e,n}\{S_e(u_e);z_e\}
 \le q_{p,\nu_e}(\sqrt{1-\gamma}),quad e=1,2,
 \tag{5}
\]
where the pivot in (5) uses the sample covariance with local coordinate \(z_e\).  The realized rescaled covariance region is exactly \(\mathcal A_N(Y_N)\).

Let \(g(A)=\operatorname{tr}(A)-\log\det(A)-p\).  Taylor expansion at \(I_p\) yields
\[
 g(I_p+H)=\tfrac12\operatorname{tr}(H^2)+O(\lVert H\rVert^3).
 \tag{6}
\]
Positive definiteness of \(s_0\), \(\nu_e/N\to\lambda_e\), and (6) give, for every compact set \(K\) of \((u_e,z_e)\),
\[
 \sup_{(u_e,z_e)\in K}
 \left|T_{e,n}\{S_e(u_e);z_e\}
 -(z_e-u_e)^{\mathsf T}V_{0,e}^{-1}(z_e-u_e)
 \right|\longrightarrow0,
 \tag{7}
\]
where \(V_{0,1}=\lambda^{-1}V(\Sigma_{0,1})\) and \(V_{0,2}=(1-\lambda)^{-1}V(\Sigma_{0,2})\).  In half-vector coordinates, the inverse of \(V(\Sigma_{0,e})\) represents the quadratic form
\[
 H\longmapsto\tfrac12
 \operatorname{tr}(\Sigma_{0,e}^{-1}H\Sigma_{0,e}^{-1}H),
 \tag{8}
\]
which accounts for the factor \(1/2\) in (6).

At the true covariance, \(\text{lem:wishart-pivot}\), the Gaussian covariance central limit calculation (1)--(3), and (6) show that the distribution whose quantile defines \(q_{p,\nu_e}(a)\) converges to \(\chi_m^2\).  The \(\chi_m^2\) distribution function is continuous and strictly increasing at its \(a_\gamma:=\sqrt{1-\gamma}\) quantile.  Quantile convergence therefore gives
\[
 q_{p,\nu_e}(a_\gamma)\longrightarrow
 q_\infty:=q_{\chi_m^2}(a_\gamma)>0.
 \tag{9}
\]
Define
\[
 E_\gamma=prod_{e=1}^2
 \{v_e:v_e^{\mathsf T}V_{0,e}^{-1}v_e\le q_\infty\},
 \qquad
 \mathcal A_\gamma(z)=z+E_\gamma.
 \tag{10}
\]
By (4) and (9), \(E_\gamma\) is compact, full-dimensional, semialgebraic, and regular closed.

Apply \(\text{lem:skorokhod-representation}\) to (3) and work temporarily with versions for which \(Y_N\to y\) almost surely.  Fix one such path.  If \(u_N\in\mathcal A_N(Y_N)\) and \(u_N\to u\), then (7)--(9) imply
\[
 (y_e-u_e)^{\mathsf T}V_{0,e}^{-1}(y_e-u_e)\le q_\infty,
 \quad e=1,2,
 \qquad\text{so}\qquad
 \limsup_N\mathcal A_N(Y_N)\subseteq\mathcal A_\gamma(y).
 \tag{11}
\]
If \(K\subset\operatorname{int}\mathcal A_\gamma(y)\) is compact, continuity and compactness provide \(\delta>0\) for which both limiting quadratics in (7) are at most \(q_\infty-2\delta\) on \(K\).  Equations (7), (9), and \(Y_N\to y\) then give the erosion property
\[
 K\subseteq\mathcal A_N(Y_N)
 \quad\text{for all sufficiently large }N.
 \tag{12}
\]

The exact bodies also have a common compact bound.  Indeed, let \(a_{e1},\ldots,a_{ep}\) be the eigenvalues of \(S_e^{-1/2}\widehat\Sigma_eS_e^{-1/2}\).  The scalar function \(r(a)=a-\log a-1\) is nonnegative, vanishes only at one, and is bounded away from zero off each neighborhood of one.  By (5) and (9),
\[
 0\le\sum_{j=1}^p r(a_{ej})=O(N^{-1}).
 \tag{13}
\]
Taylor's inequality for \(r\) near one now yields \(\max_j|a_{ej}-1|=O(N^{-1/2})\), uniformly over the pivot sublevel.  Since \(Y_N\) is bounded on the fixed path and \(\Sigma_{0,e}\succ0\), the resulting Loewner inequalities between \(S_e\) and \(\widehat\Sigma_e\) imply
\[
 \sup_{u\in\mathcal A_N(Y_N)}\lVert u-Y_N\rVert=O(1).
 \tag{14}
\]
Thus (11), (12), and (14) are precisely the upper-limit, erosion, and common-compact-bound conditions in \(\text{lem:generic-ellipsoid-intersection-stability}\).

For \(b\in\mathbb B(s_0)\), define
\[
 C_{b,N}=\sqrt N(\overline{\mathcal S_b}-s_0),
 \qquad
 Q_{b,N}=\overline{C_{b,N}\cap\mathcal A_N(Y_N)},
 \qquad
 Q_b(y)=T_b(s_0)\cap\mathcal A_\gamma(y).
 \tag{15}
\]
The condition \(\text{ass:tangent-cone-regularity}\), evaluated at \(t=N^{-1/2}\), states
\[
 C_{b,N}\longrightarrow T_b(s_0)
 \tag{16}
\]
locally in the Painlev\'e--Kuratowski topology.  Each cone \(T_b(s_0)\) is semialgebraic: the scaled graph used in the definition of the tangent cone is semialgebraic after introducing the positive scale variable, closure preserves semialgebraicity, and projection preserves it by \(\text{lem:tarski-seidenberg-projection}\).  There are finitely many accessible branches by \(\text{thm:finite-branch-representation}\).  Hence the union, over accessible \(b\), of the exceptional translation sets in \(\text{lem:generic-ellipsoid-intersection-stability}\) is Lebesgue null.  Because \(y\) has a density by (4), almost surely every accessible branch is simultaneously at a stable translation.

On this probability-one event, (11)--(16) and \(\text{lem:generic-ellipsoid-intersection-stability}\) give
\[
 Q_{b,N}\longrightarrow Q_b(y)
 \tag{17}
\]
in the Fell topology, with Hausdorff convergence if \(Q_b(y)\ne\varnothing\), and with \(Q_{b,N}=\varnothing\) eventually if \(Q_b(y)=\varnothing\).  The use of the closed cell in (15) does not enlarge the eventual branch image.  Indeed, \(\mathcal S_b\) is dense in \(\overline{\mathcal S_b}\).  In the lower-limit proof of the intersection lemma, a limiting point is first approximated inside \(\operatorname{int}\mathcal A_\gamma(y)\); a closed ball around that point is contained in \(\mathcal A_N(Y_N)\) eventually by (12).  Any approximating point from \(C_{b,N}\) can therefore be perturbed within that ball into \(\sqrt N(\mathcal S_b-s_0)\).  The upper limit is unchanged because \(\mathcal S_b\subseteq\overline{\mathcal S_b}\).  Consequently (17) also holds when the left side is the closure of the rescaled intersection with \(\mathcal S_b\) itself.

If \(b\notin\mathbb B(s_0)\), then \(s_0\notin\overline{\mathcal S_b}\).  Since the latter set is closed, it misses a neighborhood of \(s_0\); (14) therefore makes its intersection with the covariance region empty for all sufficiently large \(N\).  This disposes of every inaccessible branch.

It remains to map (17) through the causal target maps.  By \(\text{thm:finite-branch-representation}\) and \(\text{def:confidence-union}\), the confidence union is the closure of the finite union of the images \(\theta_b\) over the corresponding covariance intersections.  The condition \(\text{ass:branch-smoothness}\), together with the common compact bound (14), gives uniformly over the relevant \(u\),
\[
 \theta_b(s_0+u/\sqrt N)=\theta_b(s_0)+o(1).
 \tag{18}
\]
For a nonempty limiting intersection, (17)--(18) give convergence of the branch image to \(\{\theta_b(s_0)\}\); for an empty limiting intersection, the branch is eventually absent.  Taking the finite union yields, on the Skorokhod coupling,
\[
 \widehat{\mathcal C}_{n,\gamma}\longrightarrow
 \{\theta_b(s_0):b\in\mathbb B(s_0),\ Q_b(y)\ne\varnothing\}
 =\mathcal K_h^{(0)}
 \tag{19}
\]
in the Fell topology.  Almost-sure convergence of the coupled versions implies the first asserted convergence in distribution under the original triangular laws.

Finally suppose \(\Gamma(s_0)=\{\tau_0\}\).  By the definition of \(\Gamma(s_0)\), \(\theta_b(s_0)=\tau_0\) for every accessible \(b\).  Continuous differentiability and (14) sharpen (18) to the uniform expansion
\[
 \sqrt N\{\theta_b(s_0+u/\sqrt N)-\tau_0\}
 =D\theta_b(s_0)u+o(1).
 \tag{20}
\]
Apply (17) and (20) branch by branch and again take the finite union.  This gives
\[
 \sqrt N(\widehat{\mathcal C}_{n,\gamma}-\tau_0)
 \longrightarrow
 \bigcup_{b\in\mathbb B(s_0)}
 \{D\theta_b(s_0)u:u\in T_b(s_0)\cap
 \mathcal A_\gamma(h+\mathbb Z_0)\}
 =\mathcal K_h^{(1)}
 \tag{21}
\]
almost surely on the coupling and hence in distribution, in the Fell topology, under the original laws.  Empty limiting intersections are permitted, which is why Fell convergence, rather than only convergence on the nonempty Hausdorff hyperspace, is asserted.